\chapter{Future work}
\label{ch:FutureWork}

Both parts of the approach leave room for further work.
Most of the directions below start from the drawbacks that \autoref{ch:Conclusion} names.

\section{Feature annotations}
\label{sec:FutureWork:FeatureAnnotations}

Validating a configuration prior to derivation would eliminate the need for the methodologist to manually ensure its validity.
The preprocessor derives every configuration provided, so selecting two alternatives of \textit{ClassCreation} yields two reactions on the same trigger instead of an error.
The feature model for the umljava rules is already available as \texttt{feature-model.uvl}, located alongside the nine configurations \cite{vitruv-dsls-fork}.
A preprocessor capable of interpreting this model could reject invalid configurations and explicitly identify the violated group or constraint.
White et al. identify the minimal modification to a selection required to transform an invalid configuration into a valid one \cite{white2010automated}.
Validity corresponds to satisfiability, a property that existing feature model tools are already capable of verifying \cite{flamapy-ide,kastner2009featureide}.
Feature names used in annotations and configurations could be validated against the same model, as misspelled names currently result in dropped reactions without any warning.

A runtime decision, such as the collection type dialog described in \autoref{listing:existing-user-interaction}, could be integrated into the configuration.
The options \texttt{ArrayList}, \texttt{LinkedList}, and \texttt{HashSet} would constitute an alternative group, with the selection determining the decision at derivation time.
However, the dialog is initiated by a routine accessed by the reactions of several features, including \textit{AttributeType}, \textit{Multiplicity}, and \textit{ReturnType}.
Because routines lack annotations, each reaction would require a separate copy for each collection type.
Applying annotations to routines would eliminate the need for these duplicate copies.
In this approach, shared rules would invoke different routines, and switching configurations would necessitate comparing semantic hashes.
If choices differ between elements, the dialog remains necessary, since a configuration applies a single decision to the entire rule set (\autoref{sec:Concept:FeatureAnnotations:BindingTime}).

\section{Migration}
\label{sec:FutureWork:Migration}

Combining the two commits involved in a selective migration can decrease overall runtime.
A full migration replays its snapshot in a single commit, while a selective migration requires two separate commits.
The first commit removes the affected elements, and the second commit reinserts their respective copies.
Each commit constitutes a complete change propagation, concluding with writing the modified models, correspondences, identifiers, and the registry to disk.
A single commit would accomplish the same tasks within one propagation and write the resulting state only once.
However, Vitruvius records all deletions at the conclusion of each commit.
Therefore, a single commit must process deletions prior to insertions to ensure that old counterparts are removed before new ones are established.
Writing the state only once at the conclusion of the migration would also reduce the computational cost of subsequent commits, including those resulting from source updates.
In this scenario, the preservation step would be required to read the migrated state from memory rather than from disk.

Utilizing copies of the affected elements may reduce the duration of the preservation step in selective migration.
Currently, this step requires approximately the same amount of time as a full migration (\autoref{sec:Evaluation:MigrationEffort:Runtime}).
The process involves loading both the pre-migration and post-migration states from disk, then comparing every derived file, including those that remain unchanged.
Selective migration already duplicates the affected elements in memory prior to the initial commit.
If extended to include the derived elements corresponding to the affected subtrees, these copies would also contain the underived content of the dismantled elements.
Consequently, the preservation step could begin with these in-memory copies rather than loading both states from disk.

Maintaining a persistent record of the changes performed by each rule would enhance the precision of the preservation step.
The preservation step identifies underived content through the absence of a corresponding element.
In addition to the handwritten content, the reports list 957 items that were not authored by any individual, primarily consisting of modifiers, initial values, and methods (\autoref{sec:Evaluation:MigrationQuality:ContentNotCarriedOver}).
During a propagation, Vitruvius separates the change a developer commits from the consequential changes of the rules, but it does not persist this separation.
Documenting the specific rule responsible for each change would enable the preservation step to exclude all content generated by rules that has not been subsequently modified by a developer.
Hearnden et al. maintain a comparable record for a continuously running transformation \cite{hearnden2006incremental}, whereas a V-SUM would require persistence of this record between executions.

The same record enables the removal of residual elements from the dominant model.
Five features also modify the UML model, and the residual elements from these features cause 43 out of 72 selective migrations to deviate from derivations performed from scratch (\autoref{sec:Evaluation:MigrationQuality:EquivalenceToADerivationFromScratch}).
A full migration also retains the dominant model and is subject to the same limitation.
Both types of migration can subsequently reverse the recorded changes associated with deselected features.

Given such a record, the preservation process can automatically restore handwritten content.
In the 27 migrations where the report indicates that the complete body of \texttt{Member.totalWeight} is retained, the \texttt{--preserve user} option does not add items from the previous rules.
Content identified as lost in the report can be recovered through targeted repairs, each addressing a specific recorded cause.
In the eight migrations to config7, \texttt{Member} is transformed into an interface, with methods that contain no statements.
A repair may involve creating a class that implements the interface and contains the method bodies.
This approach is similar to the delegation mechanism, where a handwritten class implements an interface generated from the model's method signatures \cite{greifenberg2015integration}.
Because the new rules propagate the created class to other models, any repair must be compatible with the target configuration.
Once handwritten content is restored and residual changes are reverted, migrating to another configuration and returning can reproduce the original V-SUM, except for the order of modifiers and members.
